\centerline{\bf \Huge Алгебра EZ}

\begin{flushright}
        \scriptsize{Герман, переходи в тех лицей!!!}
\end{flushright}

\problem{1}
Сумма четырех целых чисел равна 0. Числа расставили по кругу и каждое умножили на сумму двух его соседей. Докажите, что сумма этих четырех произведений, умноженная на $-1$ , равна удвоенному квадрату целого числа.

\problem{2} 
Герман выбрал вещественные числа $a, b, c_1, c_2, \ldots, c_{100}$ и выписал на доску $100$ квадратных трехчленов $f_i(x)=a x^2+b x+c_i$. Оказалось, что все они имеют два вещественных корня. Для трехчлена $f_i(x)$ обозначим через $x_i$ один из его корней. Какие значения может принимать сумма

$$
f_2\left(x_1\right)+f_3\left(x_2\right)+\ldots f_{100}\left(x_{99}\right)+f_1\left(x_{100}\right)?
$$

\problem{3}
Даны действительные числа $a, b$ и $c$, не все из которых равны между собой. Докажите, что равенство $a+b+c=0$ выполнено тогда и только тогда, когда $a^2+a b+b^2=b^2+b c+c^2=c^2+c a+a^2$.

\problem{4}
Даны четыре последовательных натуральных числа, больших 100. Докажите, что из них можно выбрать три числа, сумма которых представляется в виде произведения трёх различных натуральных чисел, больших 1.

\problem{5}
Числа $x, y, z$ таковы, что $x y z=1$. Вычислите значение выражения

$$
\frac{1}{1+x+x y}+\frac{1}{1+y+y z}+\frac{1}{1+z+x z}
$$


\problem{6}
Даны действительные числа $a$ и $b$, причём $b>a>1$. Пусть

$$
x_n=2^n(\sqrt[2^n]{b}-\sqrt[2^n]{a})
$$


Докажите, что последовательность $x_1, x_2, \ldots$ убывает.

\problem{7}
Уравнение с целыми коэффициентами $x^4+a x^3+b x^2+c x+d=0$ имеет четыре положительных корня с учетом кратности. Найдите наименьшее возможное значение коэффициента $b$ при этих условиях.

\newpage

\hspace{3mm}

\centerline{\bf \Huge Алгебра Hard}

\begin{flushright}
        \scriptsize{А то на прогиб кину.}
\end{flushright}

\problem{1}
Обозначим через $f(x)$ функцию, которая равна 1 при любом целом $x$ и равна 0 при остальных $x$. Елена Михайловна дала задание двоечнику Глебу записать функцию $f(x)$ с помощью букв $x$, целых чисел, знаков сложения, вычитания, умножения, деления и операции взятия целой части. Помогите Глебу.

\problem{2}
Для положительных чисел $x_1, \ldots, x_n$ докажите неравенство:

$$
\left(\frac{x_1}{x_2}\right)^4+\left(\frac{x_2}{x_3}\right)^4+\ldots+\left(\frac{x_n}{x_1}\right)^4 \geqslant \frac{x_1}{x_5}+\frac{x_2}{x_6}+\ldots+\frac{x_{n-3}}{x_1}+\frac{x_{n-2}}{x_2}+\frac{x_{n-1}}{x_3}+\frac{x_n}{x_4}
$$

\problem{3}
Выражение $x^4+x^3-3 x^2+x+2$ возвели в натуральную степень, раскрыли скобки и привели подобные слагаемые. Докажите, что хотя бы один коэффициент полученного выражения будет отрицательным.

 
\problem{4}
Числа $x$ и $y$ удовлетворяют условиям $x^2+y^2=1$ и $20 x^3-15 x=3$. Найдите $\left|20 y^3-15 y\right|$.

\problem{5}
Многочлен $P(x)$ удовлетворяет следующим условиям: $P(0)=1,(P(x))^2=1+x+x^{100} Q(x)$ при всех действительных $x$, где $Q(x)$ - некоторый многочлен. Докажите, что коэффициент при $x^{99}$ многочлена $(P(x)+1)^{100}$ равен нулю.